\abstract{%
Reported gains from offloading, swapping, caching, and compression in
hierarchical-memory inference are hard to compare, partly because the operating
point goes unreported and common bandwidth ratios are circular, placing total
step time in both metric and outcome. We propose a
measurement discipline: describe any operating point with two independently
measurable, scale-invariant ratios---residency $R_C=C_{\mathrm{fast}}/W$ and
non-circular overlap $R_B=T_{\mathrm{comp}}/T_{\mathrm{transfer}}$---plus the
fields that make them recomputable. A conservative service-time bound
yields two definitional floors: non-resident storage $\max(0,1-R_C)$ and exposed
transfer $\max(0,1-R_B)$. The same audit applies to weights, key--value caches,
embeddings, and retrieval indices. Proof-of-concept probes on one CPU and
\NumAccel{} accelerators of two vendors show the protocol is executable and
directionally consistent (\AccelCapRatioMin--\AccelCapRatioMax$\times$ under low
residency). Production effect sizes, per-device boundaries, and population-level
generality remain open and motivate a preregistered test.
}

\keywords{hierarchical memory systems, memory orchestration, AI inference,
reporting contract, working-set residency, compute--transfer overlap,
reproducible measurement}
